\chapter{Quản trị dữ liệu và quy trình CI/CD}
\label{ch:data_gov}

\section{Quản trị Dữ liệu (Data Governance) và Chất lượng Dữ liệu (Data Quality)}

Trong kiến trúc trục tích hợp dữ liệu của trường đại học, tính đa dạng của các nguồn dữ liệu đặt ra thách thức lớn về tính nhất quán. Do đó, việc thiết lập cơ chế quản trị dữ liệu (Data Governance) song hành với các giải pháp kỹ thuật là yêu cầu bắt buộc.

Hệ thống áp dụng ba trụ cột quản trị chính:

\subsection{Quản trị Dữ liệu (Data Governance)}

\subsubsection{Quản lý Siêu dữ liệu (Metadata Management)}
Để đảm bảo sự thấu hiểu chung về dữ liệu giữa các phòng ban, hệ thống xây dựng bộ từ điển dữ liệu (Data Dictionary) chuẩn. Metadata bao gồm định nghĩa trường, kiểu dữ liệu, các danh mục mã (code list) và mối quan hệ thực thể. Metadata đóng vai trò là "ngôn ngữ chung", giúp giảm thiểu xung đột ngữ nghĩa khi tích hợp.

\subsubsection{Quyền sở hữu dữ liệu (Data Ownership \& Stewardship)}
Trách nhiệm về dữ liệu được phân định rõ ràng theo từng miền (domain):
\begin{itemize}
    \item \textbf{Data Owner:} Chịu trách nhiệm pháp lý và phê duyệt các chính sách chất lượng dữ liệu.
    \item \textbf{Data Steward:} Chịu trách nhiệm vận hành, rà soát và xử lý các vấn đề dữ liệu ở mức kỹ thuật.
\end{itemize}

\subsubsection{Kiểm soát truy cập (Data Access Control)}
Hệ thống tuân thủ nguyên tắc đặc quyền tối thiểu (Least Privilege). Tầng cơ sở dữ liệu áp dụng các chính sách phân quyền (\texttt{GRANT}/\texttt{REVOKE}), trong khi tầng ứng dụng và API Gateway kiểm soát quyền truy cập thông qua cơ chế RBAC và xác thực bằng JWT.

\subsection{Kiểm soát Chất lượng Dữ liệu (Data Quality - DQ)}

Mọi luồng dữ liệu khi đi qua Data Integration Server đều phải vượt qua các cổng kiểm soát chất lượng (Quality Gates) nghiêm ngặt, đảm bảo tính nhất quán theo tiêu chuẩn \cite{gartner:data_consistency}:

\begin{itemize}
    \item \textbf{Tính hợp lệ (Validity Checks):} Kiểm tra định dạng (Regex cho email, MSSV), ràng buộc miền giá trị và tính toàn vẹn (null/not-null).
    \item \textbf{Tính nhất quán (Consistency Checks):} Đảm bảo dữ liệu quan hệ khớp với metadata (ví dụ: Mã Khoa trong hồ sơ sinh viên phải tồn tại trong danh mục Khoa).
    \item \textbf{Tính kịp thời (Timeliness Checks):} Phát hiện và đánh dấu các bản ghi trễ hạn (late-arriving data) để có quy trình xử lý riêng.
    \item \textbf{Phát hiện trùng lặp (Duplicate Detection):} Sử dụng thuật toán so khớp trên các trường định danh (CCCD, Email, MSSV) để loại bỏ dữ liệu thừa.
\end{itemize}

\subsection{Cơ chế xử lý lỗi dữ liệu}

Dữ liệu không đạt chuẩn sẽ không bị loại bỏ hoàn toàn mà được chuyển hướng vào hàng đợi lỗi (\textit{Error Queue/Dead Letter Queue}). Sự kiện này được ghi log chi tiết vào Elasticsearch và gửi cảnh báo đến Data Steward để xử lý thủ công, đảm bảo không làm gián đoạn luồng dữ liệu chính.

\section{CI/CD Tự Động Với GitHub Actions}
\label{sec:cicd-github}

\subsection{Đặt vấn đề}

Trong giai đoạn phát triển, nhóm thực hiện deploy thủ công lên EC2 qua SSH
sau mỗi lần thay đổi mã nguồn. Quy trình này tốn thời gian, dễ bỏ sót bước
và không đảm bảo tính nhất quán giữa môi trường phát triển và môi trường
production. CI/CD (Continuous Integration / Continuous Deployment) tự động
hóa toàn bộ quy trình: từ khi developer \texttt{git push} lên nhánh
\texttt{master} cho đến khi phiên bản mới chạy trên EC2 --- không cần thao
tác thủ công.

\subsection{Kiến trúc pipeline}

Pipeline CI/CD được triển khai qua \textbf{GitHub Actions},
kích hoạt mỗi khi có commit được push lên nhánh \texttt{master}.
Pipeline gồm hai job chạy tuần tự:

\begin{enumerate}
  \item \textbf{Build \& Push}: Build Docker image cho backend (NestJS)
        và frontend (Next.js), đẩy lên Docker Hub với tag \texttt{latest}.
  \item \textbf{Deploy}: SSH vào EC2 (Amazon Linux~2023, t3.medium),
        pull image mới và khởi động lại container bằng Docker~Compose.
\end{enumerate}

\begin{figure}[H]
\centering
\begin{tikzpicture}[
  node distance=1.6cm,
  box/.style={rectangle, rounded corners=4pt,
              draw=blue!55, fill=blue!8,
              minimum width=2.6cm, minimum height=0.8cm,
              font=\small\bfseries, align=center, text=blue!80},
  arr/.style={-{Stealth[length=5pt,width=4pt]}, thick, black!55}
]
  \node[box] (push)   {git push\\master};
  \node[box, right=1.4cm of push]  (trig)   {GitHub\\Actions};
  \node[box, right=1.4cm of trig]  (build)  {Build\\Docker Image};
  \node[box, right=1.4cm of build] (hub)    {Push\\Docker Hub};
  \node[box, right=1.4cm of hub]   (deploy) {SSH Deploy\\EC2};

  \draw[arr] (push)  -- (trig);
  \draw[arr] (trig)  -- (build);
  \draw[arr] (build) -- (hub);
  \draw[arr] (hub)   -- (deploy);
\end{tikzpicture}
\caption{Luồng CI/CD từ push code đến deploy EC2}
\label{fig:cicd_pipeline}
\end{figure}

\subsection{File workflow GitHub Actions}

File workflow được đặt tại \texttt{.github/workflows/deploy.yml} trong
repository. Toàn bộ thông tin nhạy cảm (địa chỉ EC2, khóa SSH, thông tin
Docker Hub) được lưu trong \textbf{GitHub Secrets} thay vì hard-code.

\begin{lstlisting}[language=yaml,
                   caption={.github/workflows/deploy.yml},
                   label={lst:cicd-workflow}]
name: Build and Deploy to EC2

on:
  push:
    branches: [ master ]

env:
  IMAGE_BACKEND:  ${{ secrets.DOCKER_USERNAME }}/hub-backend
  IMAGE_FRONTEND: ${{ secrets.DOCKER_USERNAME }}/hub-frontend

jobs:
  # Job 1: Build Docker images va push len Docker Hub
  build:
    name: Build & Push Docker Images
    runs-on: ubuntu-latest

    steps:
      - name: Checkout source
        uses: actions/checkout@v4

      - name: Log in to Docker Hub
        uses: docker/login-action@v3
        with:
          username: ${{ secrets.DOCKER_USERNAME }}
          password: ${{ secrets.DOCKER_PASSWORD }}

      - name: Build & push Backend image
        uses: docker/build-push-action@v5
        with:
          context: ./backend
          push: true
          tags: ${{ env.IMAGE_BACKEND }}:latest

      - name: Build & push Frontend image
        uses: docker/build-push-action@v5
        with:
          context: ./frontend
          push: true
          tags: ${{ env.IMAGE_FRONTEND }}:latest
          build-args: |
            NEXT_PUBLIC_KONG_URL=${{ secrets.NEXT_PUBLIC_KONG_URL }}

  # Job 2: SSH vao EC2, pull image moi va restart container
  deploy:
    name: Deploy to EC2
    runs-on: ubuntu-latest
    needs: build

    steps:
      - name: SSH and deploy
        uses: appleboy/ssh-action@v1.0.3
        with:
          host:     ${{ secrets.EC2_HOST }}
          username: ${{ secrets.EC2_USERNAME }}
          key:      ${{ secrets.EC2_SSH_KEY }}
          script: |
            cd ~/Data_Integration_Hub

            # Pull image moi tu Docker Hub
            docker compose pull backend frontend

            # Restart chi hai service thay doi
            docker compose up -d --no-build backend frontend

            # Xoa image cu de tiet kiem disk
            docker image prune -f

            echo "Deploy hoan thanh: $(date)"
\end{lstlisting}

\subsection{Cấu hình GitHub Secrets}

Trước khi workflow có thể chạy, cần khai báo các secret sau tại
\textbf{Settings $\to$ Secrets and variables $\to$ Actions} của repository:

\begin{table}[H]
\centering
\caption{Danh sách GitHub Secrets cần cấu hình}
\label{tab:github-secrets}
\begin{tabular}{lll}
\hline
\textbf{Secret} & \textbf{Giá trị} & \textbf{Mục đích} \\
\hline
\texttt{DOCKER\_USERNAME}        & Tên tài khoản Docker Hub    & Đăng nhập Docker Hub \\
\texttt{DOCKER\_PASSWORD}        & Access Token Docker Hub     & Xác thực push image \\
\texttt{EC2\_HOST}               & Elastic IP của EC2          & Địa chỉ SSH \\
\texttt{EC2\_USERNAME}           & \texttt{ec2-user}           & Tài khoản SSH \\
\texttt{EC2\_SSH\_KEY}           & Nội dung file \texttt{.pem} & Khóa xác thực SSH \\
\texttt{NEXT\_PUBLIC\_KONG\_URL} & \texttt{http://<IP>:8000}   & Build-time env Next.js \\
\hline
\end{tabular}
\end{table}

\subsection{Lưu ý kỹ thuật}

\paragraph{NEXT\_PUBLIC\_KONG\_URL phải khai báo lúc build.}
Next.js nhúng (\textit{bake}) giá trị biến môi trường
\texttt{NEXT\_PUBLIC\_*} vào bundle JavaScript tại thời điểm build,
không phải lúc runtime. Biến này phải được truyền qua \texttt{build-args}
trong bước build Frontend --- không thể truyền vào container sau khi đã
build xong.

\paragraph{Chỉ pull backend và frontend, không restart toàn bộ stack.}
Các service nền tảng như PostgreSQL, MinIO, Kong không thay đổi qua mỗi
lần deploy. Lệnh \texttt{docker compose pull backend frontend} chỉ kéo
image của hai service thay đổi, giúp giảm thời gian deploy từ 10--15 phút
xuống còn khoảng 2--3 phút.

\paragraph{Elastic IP là bắt buộc.}
EC2 cấp phát Public IP động --- địa chỉ thay đổi mỗi khi instance
Stop/Start. Nếu không gán Elastic IP tĩnh, secret \texttt{EC2\_HOST} sẽ
bị sai sau lần khởi động lại đầu tiên và toàn bộ pipeline thất bại ở bước
deploy.

\subsection{Kiểm thử tự động (Unit Testing)}

Song song với việc thiết lập pipeline CI/CD, nhóm xây dựng bộ unit test
cho các thành phần cốt lõi của backend bằng \textbf{Jest} --- framework
kiểm thử mặc định của NestJS. Tổng cộng \textbf{41 test case} được viết
trên 3 file, bao phủ các logic nghiệp vụ quan trọng nhất:

\begin{table}[H]
\centering
\caption{Bộ unit test của backend}
\label{tab:unit-tests}
\begin{tabular}{llp{7cm}}
\hline
\textbf{File} & \textbf{Tests} & \textbf{Nội dung kiểm thử} \\
\hline
\texttt{string.util.spec.ts}           & 6  & Hàm \texttt{snakeToCamel} chuyển đổi tên bảng \\
\texttt{sync-engine.service.spec.ts}   & 20 & Dedup records, phát hiện field lạ, phát hiện orphan, xóa theo ID \\
\texttt{data-integration.service.spec.ts} & 15 & Lọc records theo thời gian, cơ chế retry, xử lý pipeline rỗng \\
\hline
\end{tabular}
\end{table}

Toàn bộ test được chạy bằng lệnh \texttt{npm test} và pass 100\%.
Bước tiếp theo là tích hợp bộ test này vào pipeline GitHub Actions:
thêm job \texttt{test} chạy \textbf{trước} job \texttt{build}, đảm bảo
mọi commit lên \texttt{master} đều phải vượt qua kiểm thử tự động trước
khi Docker image được build và deploy lên production.

\subsection{Đánh giá}

\textbf{Ưu điểm:}
\begin{itemize}
  \item \textbf{Tự động hóa hoàn toàn}: Developer chỉ cần \texttt{git push},
        phiên bản mới lên production sau 3--5 phút.
  \item \textbf{Tách biệt build và runtime}: Image được build trên
        GitHub-hosted runner (Ubuntu 24.04), không tiêu tốn CPU/RAM của EC2.
  \item \textbf{Rollback dễ dàng}: Docker Hub lưu các image theo tag.
        Rollback bằng cách chỉnh tag trong \texttt{docker-compose.yml}
        và chạy lại \texttt{docker compose up -d}.
  \item \textbf{Nền tảng kiểm thử sẵn sàng}: Bộ 41 unit test đã được
        xây dựng và pass 100\%, sẵn sàng tích hợp vào pipeline.
\end{itemize}

\textbf{Hướng mở rộng:}
Tích hợp job \texttt{test} vào đầu pipeline để mọi commit được kiểm thử
tự động trước khi build. Ngoài ra có thể bổ sung \texttt{npm run test:cov}
để theo dõi độ phủ test (code coverage) qua từng lần merge.

\chapter{Chiến lược Sao lưu và Phục hồi Dữ liệu}
\label{ch:backup_strategy}

Nhóm thực hiện đề xuất giải pháp sao lưu sử dụng \textbf{pgBackRest} kết hợp với hệ thống lưu trữ đối tượng \textbf{MinIO} triển khai nội bộ. Giải pháp này đảm bảo tính bảo mật dữ liệu của trường đại học trong khi vẫn tận dụng được sự linh hoạt của kiến trúc S3.

\begin{table}[H]
\centering
\caption{So sánh MinIO với các giải pháp lưu trữ đám mây}
\label{tab:backup_comparison}
\begin{tabular}{|p{3cm}|p{3.5cm}|p{3cm}|p{3cm}|}
\hline
\textbf{Tiêu chí} & \textbf{MinIO (On-premise)} & \textbf{AWS S3} & \textbf{Google Cloud Storage} \\
\hline
Mô hình triển khai & Private Cloud (Nội bộ) & Public Cloud & Public Cloud \\
\hline
Vị trí dữ liệu & Mạng nội bộ nhà trường & Datacenter nhà cung cấp & Datacenter nhà cung cấp \\
\hline
Khả năng tích hợp & Tương thích hoàn toàn S3 API & Native & SDK riêng \\
\hline
Chi phí vận hành & Thấp (Tận dụng phần cứng sẵn có) & Cao (Theo dung lượng) & Cao (Theo dung lượng) \\
\hline
Độ trễ & Rất thấp (LAN) & Trung bình & Trung bình \\
\hline
\textbf{Độ phù hợp} & \textbf{Rất cao} & Thấp & Thấp \\
\hline
\end{tabular}
\end{table}

\section{Chiến lược và Chu kỳ sao lưu}

Để cân bằng giữa hiệu năng hệ thống và an toàn dữ liệu, hệ thống áp dụng chiến lược sao lưu hỗn hợp:

\begin{itemize}
    \item \textbf{Full Backup (Toàn phần):} Thực hiện \textbf{1 lần/tuần} vào thời điểm thấp điểm (00:00 – 03:00 Chủ nhật). Bản sao lưu này chứa toàn bộ dữ liệu hệ thống.
    \item \textbf{Incremental Backup (Gia tăng):} Thực hiện \textbf{hằng ngày} vào lúc 01:00 sáng. Bản sao lưu này chỉ chứa các dữ liệu thay đổi so với lần backup trước đó, giúp tiết kiệm dung lượng và thời gian.
    \item \textbf{Lưu trữ:} Các bản sao lưu được giữ lại ít nhất 1 tháng (Retention Policy).
\end{itemize}

\section{Phương thức kỹ thuật}

Hệ thống sử dụng công cụ \textbf{pgBackRest} nhờ các ưu điểm vượt trội về hiệu năng và khả năng tích hợp S3. Quy trình xử lý file backup trước khi đẩy lên MinIO bao gồm:
\begin{enumerate}
    \item \textbf{Nén dữ liệu:} Sử dụng thuật toán Gzip/LZ4 để giảm dung lượng đường truyền và lưu trữ.
    \item \textbf{Mã hóa:} Áp dụng mã hóa AES-256 để đảm bảo dữ liệu không thể bị đọc trộm ngay cả khi hệ thống lưu trữ bị xâm nhập.
\end{enumerate}

\section{Cam kết chỉ số phục hồi (RPO \& RTO)}

Dựa trên chiến lược sao lưu toàn phần kết hợp sao lưu gia tăng hằng ngày, hệ thống cam kết:

\begin{itemize}
    \item \textbf{Recovery Point Objective (RPO): 1 ngày.}
    Trong trường hợp xấu nhất, lượng dữ liệu tối đa bị mất tương đương với các giao dịch phát sinh từ sau lần Incremental Backup gần nhất (tối đa 24 giờ).
    
    \item \textbf{Recovery Time Objective (RTO): 1–3 giờ.}
    Thời gian cần thiết để tải các bản backup từ MinIO, giải mã, giải nén và khôi phục vào hệ quản trị cơ sở dữ liệu.
\end{itemize}